\documentclass[10pt,a4paper]{amsart}
\usepackage[margin=19mm]{geometry}
\usepackage{amsmath,amssymb,mathtools}
\usepackage[scr=rsfs,cal=euler]{mathalfa}
\usepackage{xfrac}
\usepackage[unicode,hidelinks,bookmarks=false]{hyperref}
\newcommand{\ie}{i.e.\ }
\newcommand{\cf}{cf.\ }
\newcommand{\resp}{respectively\ }
\newcommand{\Subsection}{\subsection}
\providecommand{\nobreakdash}{\nobreak\hbox{-}}
\setlength{\emergencystretch}{2em}
\multlinegap0pt
\usepackage{bm}
\usepackage{bbm}
\usepackage{dsfont}
\usepackage{xspace}
\usepackage{cancel}
\usepackage{mathtools}
\usepackage{amsthm}
\makeatletter
\@ifundefined{lemma}{\newtheorem{lemma}{Lemma}}{}
\@ifundefined{proposition}{\newtheorem{proposition}{Proposition}}{}
\makeatother
\title[Nonexistence of weakly stable Yang-Mills fields]{Nonexistence of weakly stable Yang-Mills fields}
\author{Xiaoli Han, Yang Wen}
\date{Source: arXiv:2601.17649v1 (CC BY 4.0)}
\begin{document}
\maketitle

\noindent\emph{Source and licence.} Nonexistence of weakly stable Yang-Mills fields. Version arXiv:2601.17649v1 (CC BY 4.0), distributed under the Creative Commons Attribution 4.0 International licence, as recorded in the arXiv licence metadata for this exact version and shown on the abstract page https://arxiv.org/abs/2601.17649v1. Full rights notice: licenses/arxiv-2601.17649-CC-BY-4.0.txt.\par\smallskip

\noindent\textit{Editorial scope.} Counted unit: the Proposition whose source label is \texttt{sum L(B\_V\_k)} in the article (source0.tex lines 321--327, source label \texttt{sum L(B\_V\_k)}), delivered together with the article's own complete original proof of it (source0.tex lines 328--385, one \texttt{proof} environment of 58 lines). No mathematics is written, repaired or re-derived: statement and proof are the article's own text line for line. Required context retained uncounted: conformal vector field (lines 263--273); S(iV\_R) (lines 280--287). Every definition, display and label of those lines is retained unchanged. Omitted from this package and not redistributed: 1--262, 274--279, 288--320, 386--701. Nothing inside a retained excerpt is omitted or altered: every retained body is delivered exactly as the article's lines are written, so no omission receipt of any kind accompanies this package. Citations: the retained excerpts carry no citation command at all, so no bibliography entry of the article is transcribed and no reference is redistributed. Reference adaptation: none was needed and none was made; every reference inside the delivered unit resolves inside it, so no unresolved reference or citation is produced. Printed slips kept literal: no printed word, symbol, hypothesis or number of the counted statement or of its proof was altered. Layout and class: the article's own class is \texttt{amsart} with packages including geometry, amssymb, amsmath, mathtools, verbatim, bm, hyperref, footmisc, natbib, mathrsfs, caption, subcaption, amsthm; the wrapper is the lane's pinned \texttt{amsart} editorial wrapper at 10pt on A4, which restates the theorem-like environments the retained text needs and adds the article's own macro definitions that the retained text needs (none), with extra wrapper packages (bm, bbm, dsfont, xspace, cancel, mathtools). The delivered numbering is the unit's own. Assets: no retained span contains a figure environment, a graphics-inclusion command, a PDF-page inclusion, a table or a TikZ picture; no image file is copied into this package, the package declares no asset dependency and no image file is redistributed with it. Licence: version arXiv:2601.17649v1 of this article is distributed under the Creative Commons Attribution 4.0 International licence, as recorded in the arXiv licence metadata for this exact version and shown on the abstract page https://arxiv.org/abs/2601.17649v1. A full rights notice is delivered at licenses/arxiv-2601.17649-CC-BY-4.0.txt. Characters: every retained excerpt is pure ASCII, so the delivered engine, XeLaTeX with its default Latin Modern text font, reports no missing character.
	\begin{proposition}\label{conformal vector field}
		For any vector $v\in\mathbb{R}^{n+1}$, let $F_v(x)=v\cdot x:\mathbb{R}^{n+1}\to\mathbb{R}$ be the linear map and $f_v=F_v\mid_{S^n}$ be the restriction on $S^n$. Then $V=\operatorname{grad}(f_v)$ is a conformal vector field of $(S^n,g)$ and satisfies
		\begin{align*}
			&D_XV=-f_vX,\\
			&D^\ast DV=V
		\end{align*}
		for any $X\in\mathscr{X}(S^n)$, where
		\begin{align*}
			D^\ast DV=-\sum_{i=1}^n(D_{e_i}D_{e_i}V-D_{D_{e_i}e_i}V).
		\end{align*}
	\end{proposition}

	\begin{lemma}\label{S(iV_R)}
		For any $v\in\mathbb{R}^{n+1}$ and $X\in\mathscr{X}(S^n)$, we have
		\begin{align*}
			\mathscr{S}(B_v)(X)&=(4-n+\lambda\Delta\varphi-(\lambda^2+(n-4)\lambda)|\operatorname{grad}(\varphi)|_g^2)R^\nabla(\tilde V,X)+(4-n+3\lambda)e^{\lambda\varphi}f_vR^\nabla(\operatorname{grad}(\varphi),X)\\
			&-\lambda^2X(\varphi)R^\nabla(\operatorname{grad}(\varphi),\tilde V)-\lambda\nabla_{\operatorname{grad}(\varphi)}R^\nabla(\tilde V,X)-\lambda\nabla_{\tilde V}R^\nabla(\operatorname{grad}(\varphi),X)\\
			&+(n-4)R^\nabla(D_{\tilde V}\operatorname{grad}(\varphi),X)+\lambda R^\nabla(\tilde V,D_X\operatorname{grad}(\varphi)).
		\end{align*}
	\end{lemma}

	\begin{proposition}\label{sum L(B_V_k)}
		Assume $\{v_1,...,v_{n+1}\}$ be an orthogonal basis of $\mathbb{R}^{n+1}$, then we have
		\begin{align}
			\sum_{k=1}^{n+1}\mathscr{L}(B_{v_k})=\int_{S^n}&e^{(n-4+2\lambda)\varphi}((\frac{4-n}2\Delta\varphi+2(\lambda+\frac{n-4}2)^2|\operatorname{grad}(\varphi)|_g^2+8-2n)|R^\nabla|_g^2\\
			&+\lambda(8-2n-\lambda)|i_{\operatorname{grad}(\varphi)}R^\nabla|^2)dV_g.
		\end{align}
	\end{proposition}

	\begin{proof}
		According to Lemma \ref{S(iV_R)}, we have
		\begin{align*}
			\mathscr{L}(B_v)=(4-n)\int_{S^n}e^{(n-4)\varphi}q(v,v)dV_g,
		\end{align*}
		where at any $x\in S^n$, $q$ is a quadratic form on $\mathbb{R}^{n+1}$ defined by
		\begin{align*}
			q(v,w)=\sum_{j=1}^n&\langle (4-n+\lambda\Delta\varphi-(\lambda^2+(n-4)\lambda)|\operatorname{grad}(\varphi)|_g^2)R^\nabla(\tilde V,e_j)+(4-n+3\lambda)e^{\lambda\varphi}f_vR^\nabla(\operatorname{grad}(\varphi),e_j)\\
			&-\lambda^2e_j(\varphi)R^\nabla(\operatorname{grad}(\varphi),\tilde V)-\lambda\nabla_{\operatorname{grad}(\varphi)}R^\nabla(\tilde V,e_j)-\lambda\nabla_{\tilde V}R^\nabla(\operatorname{grad}(\varphi),e_j)\\
			&+(n-4)R^\nabla(D_{\tilde V}\operatorname{grad}(\varphi),e_j)+\lambda R^\nabla(\tilde V,D_{e_j}\operatorname{grad}(\varphi)),R^\nabla(\tilde W,e_j)\rangle
		\end{align*}
		For any $x\in S^n$, note that $\sum_{k=1}^{n+1}q(v_k,v_k)(x)$ is independent to the choice of $\{v_k\}$, we may assume $v_{n+1}=x$ and $\{v_k\mid1\le k\le n\}$ is an orthogonal basis of $T_xS^n\subset T_x\mathbb{R}^{n+1}$. From proposition \ref{conformal vector field}, we have $V_k=e_k$ for $1\le k\le n$ and $V_{n+1}=0$ at $x$. Thus we have
		\begin{align*}
			\sum_{k=1}^{n+1}\sum_{j=1}^n\langle R^\nabla(\tilde V_k,e_j),R^\nabla(\tilde V_k,e_j)\rangle=e^{2\lambda\varphi}\sum_{i,j=1}^n\langle R^\nabla(e_i,e_j),R^\nabla(e_i,e_j)\rangle=2e^{2\lambda\varphi}|R^\nabla|_g^2.
		\end{align*}
		Since $\sum_kf_{v_k}^2\equiv1$ is a constant map, we have
		\begin{align*}
			\sum_{k=1}^{n+1}f_{v_k}\tilde V_k=\frac12e^{\lambda\varphi}\operatorname{grad}(\sum_{k=1}^{n+1}f_v^2)=0,
		\end{align*}
		and thus
		\begin{align*}
			\sum_{k=1}^{n+1}\sum_{j=1}^n\langle e^{\lambda\varphi}f_vR^\nabla(\operatorname{grad}(\varphi),e_j),R^\nabla(\tilde V_k,e_j)\rangle=0.
		\end{align*}
		We also have
		\begin{align*}
			\sum_{k=1}^{n+1}\sum_{j=1}^n\langle e_j(\varphi)R^\nabla(\operatorname{grad}(\varphi),\tilde V_k),R^\nabla(\tilde V_k,e_j)\rangle=-e^{2\lambda\varphi}|i_{\operatorname{grad}(\varphi)}R^\nabla|_g^2.
		\end{align*}
		We assume $De_i=0$ at $x$. Note that
		\begin{align*}
			&\sum_{i,j=1}^n\langle\nabla_{e_i}R^\nabla(\operatorname{grad}(\varphi),e_j),R^\nabla(e_i,e_j)\rangle\\
			=&-\sum_{i,j=1}(\langle\nabla_{\operatorname{grad}(\varphi)}R^\nabla(e_j,e_i),R^\nabla(e_i,e_j)\rangle+\langle\nabla_{e_j}R^\nabla(e_i,\operatorname{grad}(\varphi)),R^\nabla(e_i,e_j)\rangle\\
			=&\operatorname{grad}(\varphi)(|R^\nabla|_g^2)-\sum_{i,j=1}^n\langle\nabla_{e_i}R^\nabla(\operatorname{grad}(\varphi),e_j),R^\nabla(e_i,e_j)\rangle.
		\end{align*}
		Thus we have
		\begin{align*}
			\sum_{k=1}^{n+1}\sum_{j=1}^n\langle \nabla_{\tilde V_k}R^\nabla(\operatorname{grad}(\varphi),e_j),R^\nabla(\tilde V_k,e_j)\rangle=e^{2\lambda\varphi}\sum_{i,j=1}^n\langle\nabla_{e_i}R^\nabla(\operatorname{grad}(\varphi),e_j),R^\nabla(e_i,e_j)\rangle=\frac12e^{2\lambda\varphi}\operatorname{grad}(\varphi)(|R^\nabla|_g^2).
		\end{align*}
		Still assuming $De_i=0$ at $x$, we have
		\begin{align*}
			&\sum_{k=1}^{n+1}\sum_{j=1}^n\langle (n-4)R^\nabla(D_{\tilde V_k}\operatorname{grad}(\varphi),e_j)+\lambda R^\nabla(\tilde V_k,D_{e_j}\operatorname{grad}(\varphi)),R^\nabla(\tilde V_k,e_j)\rangle\\
			=&(n-4+\lambda)e^{2\lambda\varphi}\sum_{i,j=1}^n\langle R^\nabla(D_{e_i}\operatorname{grad}(\varphi),e_j),R^\nabla(e_i,e_j)\rangle
		\end{align*}
		and
		\begin{align*}
			&\sum_{i,j=1}^n\langle R^\nabla(D_{e_i}\operatorname{grad}(\varphi),e_j),R^\nabla(e_i,e_j)\rangle\\
			=&\sum_{i,j=1}^ne_i(\langle R^\nabla(\operatorname{grad}(\varphi),e_j),R^\nabla(e_i,e_j)\rangle)-\langle\nabla_{e_i}R^\nabla(\operatorname{grad}(\varphi),e_j),R^\nabla(e_i,e_j)\rangle-\langle R^\nabla(\operatorname{grad}(\varphi),e_j),\nabla_{e_i}R^\nabla(e_i,e_j)\rangle\\
			=&\operatorname{div}(\sum_{j=1}^n\langle R^\nabla(\operatorname{grad}(\varphi),e_j),R^\nabla(\cdot,e_j)\rangle)-\frac12\operatorname{grad}(\varphi)(|R^\nabla|_g^2)+(n-4)|i_{\operatorname{grad}(\varphi)}R^\nabla|_g^2.
		\end{align*}
		By direct computation, we obtain
		\begin{align*}
			-\int_{S^n}e^{(n-4+2\lambda)\varphi}\operatorname{div}(\sum_{j=1}^n\langle R^\nabla(\operatorname{grad}(\varphi),e_j),R^\nabla(\cdot,e_j)\rangle)dV_g=(n-4+2\lambda)\int_{S^n}e^{(n-4+2\lambda)\varphi}|i_{\operatorname{grad}(\varphi)}R^\nabla|_g^2dV_g
		\end{align*}
		and
		\begin{align*}
			\frac12\int_{S^n}e^{(n-4+2\lambda)\varphi}\operatorname{grad}(\varphi)(|R^\nabla|_g^2)dV_g=\frac12\int_{S^n}e^{(n-4+2\lambda)\varphi}(\Delta\varphi-(n-4+2\lambda)|\operatorname{grad}(\varphi)|_g^2)|R^\nabla|_g^2dV_g.
		\end{align*}
		Thus, we finish the proof.
	\end{proof}

\end{document}
